\section{Theory}\label{sec:theory}

This section derives the model's main implications. The key innovation is to decompose foreign-currency demand into a hedge motive and a run motive. This separation makes it possible to generate monotone foreign-currency demand, endogenous public information, crisis amplification, and welfare reversal in a single framework.

\subsection{Monotone Foreign-Currency Demand}

Total foreign-currency demand is the sum of hedge demand and run demand:

Q_FX = Q_hedge + Q_run.

Hedge demand rises with expected overvaluation and falls with the effective cost of acquiring foreign currency. Run demand rises with crisis beliefs and with the weight placed on the public signal. Because hedge demand is increasing in overvaluation, the model delivers a weakly increasing demand for foreign currency as exchange-rate misalignment worsens.

\subsection{Stablecoin Demand}

Stablecoin demand rises with overvaluation because total foreign-currency demand rises and because stablecoins have lower access costs than cash. Stablecoins therefore become a larger share of foreign-currency access as the official exchange rate becomes more distorted.

\subsection{Endogenous Information}

Stablecoin usage improves the precision of public information. As the stablecoin market becomes more active, its price becomes more informative, increasing the weight that households place on the public signal. This creates a positive feedback between stablecoin adoption and common information.

\subsection{Crisis Amplification}

Stablecoins increase crisis risk through two channels. First, they lower access costs and therefore raise total foreign-currency demand. Second, they improve public information and thereby strengthen coordination. Both effects increase the probability that aggregate demand exceeds effective defense capacity.

\subsection{Welfare Reversal}

Stablecoins improve welfare in moderately overvalued states by lowering the cost of foreign-currency access and improving tradable-goods smoothing. But they reduce welfare in sufficiently fragile states because they increase coordinated runs and amplify disruption losses in crisis states. The model therefore implies a threshold level of overvaluation below which stablecoins are beneficial and above which they are harmful.
